% Part: computability
% Chapter: recursive-functions
% Section: pr-relations

\documentclass[../../../include/open-logic-section]{subfiles}

\begin{document}

\olfileid{cmp}{rec}{prr}
\olsection{Relasi Rekursif Primitif}


\begin{defn}
Relasi $R(\vec x)$ diarani rekursif primitif yen fungsi
karakteristike,
\[
\Char{R}(\vec x) = \left\{
  \begin{array}{ll}
  1 & \mbox{yen $R(\vec x)$} \\
  0 & \mbox{yen ora}
  \end{array}
\right.
\]
iku rekursif primitif.
\end{defn}

Kanthi tembung liya, nalika kita nyebut relasi rekursif
primitif $R(\vec x)$, kang dimaksud yaiku relasi kanthi
wujud $\Char{R}(\vec
x) = 1$, ing ngendi $\Char{R}$ iku fungsi rekursif primitif
kang ngasilake 1 utawa 0 kanggo sembarang input. Umpamane,
relasi $\fn{IsZero}(x)$, kang bener yen lan mung yen $x = 0$,
cocog karo fungsi $\Char{\fn{IsZero}}$ kang ditegesi
kanthi rekursi primitif:
\begin{align*}
\Char{\fn{IsZero}}(0) & = 1,\\
\Char{\fn{IsZero}}(x+1) & = 0.
\end{align*}

Cetha yen kita bisa nyusun komposisi relasi karo fungsi
rekursif primitif liyane. Mula ing ngisor iki uga rekursif
primitif:
\begin{enumerate}
\item Relasi kesamaan, $x = y$, kang ditegesi kanthi
  $\fn{IsZero}(\left|x -
  y\right|)$
% OLPL-123: Sumber nyebut "less-than" nanging lambang lan definisi
% IsZero nganggo pangurangan kang diwatesi ing nol netepake x <= y,
% yaiku kurang saka utawa padha karo.
\item Relasi kurang saka utawa padha karo, $x \leq y$, kang
  ditegesi kanthi $\fn{IsZero}(x
  \tsub y)$. Cathetan koreksi sumber OLPL-123: jeneng less-than ing sumber ora cocog karo lambang lan rumuse; relasi iki uga bener nalika rong input padha.
\end{enumerate}


\begin{prop}
  Himpunan relasi rekursif primitif katutup tumrap operasi
  Boolean, yaiku yen $P(\vec x)$ lan $Q(\vec x)$ rekursif
  primitif, mula relasi iki uga:
  \begin{enumerate}
  \item $\lnot P(\vec x)$
  \item $P(\vec x) \land Q(\vec x)$
  \item $P(\vec x) \lor Q(\vec x)$
  \item $P(\vec x) \lif Q(\vec x)$
  \end{enumerate}
\end{prop}

\begin{proof}
  Upamana $P(\vec x)$ lan $Q(\vec x)$ rekursif primitif,
  yaiku fungsi karakteristike $\Char{P}$ lan $\Char{Q}$
  rekursif primitif. Kita kudu nuduhake yen fungsi
  karakteristik saka $\lnot P(\vec x)$ lan sakbanjure
  uga rekursif primitif.
  \[
  \Char{\lnot P}(\vec x) = \begin{cases}
    0 & \text{yen $\Char{P}(\vec x) = 1$}\\
    1 & \text{yen ora}
  \end{cases}
  \]
  Kita bisa netepake $\Char{\lnot P}(\vec x)$ minangka
  $1 \tsub \Char{P}(\vec x)$.
  \[
  \Char{P \land Q}(\vec x) = \begin{cases}
    1 & \text{yen $\Char{P}(\vec x) = \Char{Q}(\vec x) = 1$}\\
    0 & \text{yen ora}
  \end{cases}
  \]
  Kita bisa netepake $\Char{P \land Q}(\vec x)$ minangka
  $\Char{P}(\vec x) \cdot
  \Char{Q}(\vec x)$ utawa minangka $\fn{min}(\Char{P}(\vec x), \Char{Q}(\vec
  x))$. Kanthi cara kang padha,
  \begin{align*}
    \Char{P \lor Q}(\vec x) & = \fn{max}(\Char{P}(\vec x), \Char{Q}(\vec x)) \text{ lan}\\
    \Char{P \lif Q}(\vec x) & = \fn{max}(1 \tsub
  \Char{P}(\vec x), \Char{Q}(\vec x)).
  \end{align*}
\end{proof}

\begin{prop}
  Himpunan relasi rekursif primitif katutup tumrap
  kuantifikasi mawa wates, yaiku yen $R(\vec x, z)$
  relasi rekursif primitif, mula relasi iki uga:
  \begin{align*}
    & \bforall{z < y}{R(\vec x, z)} \text{ lan}\\
    & \bexists{z < y}{R(\vec x, z)}.
  \end{align*}
  $\bforall{z < y}{R(\vec x, z)}$ bener kanggo $\vec x$
  lan $y$ yen lan mung yen $R(\vec x, z)$ bener kanggo
  saben~$z$ kang kurang saka~$y$; semono uga kanggo
  $\bexists{z < y}{R(\vec x, z)}$.
\end{prop}

\begin{proof}
  Miturut konvensi, $\bforall{z < 0}{R(\vec x, z)}$
  dianggep bener (amarga ora ana $z$ kang kurang
  saka~$0$) lan $\bexists{z < 0}{R(\vec x, z)}$
  dianggep luput. Kuantor universal mawa wates tumindak
  kaya asil kali winates utawa minimum kang diulang,
  yaiku yen $P(\vec x, y) \defiff \bforall{z <
  y}{R(\vec x, z)}$, mula $\Char{P}(\vec x, y)$
  bisa ditegesi kanthi
  \begin{align*}
    \Char{P}(\vec x, 0) & = 1\\
    \Char{P}(\vec x, y+1) & =
    \fn{min}(\Char{P}(\vec x, y), \Char{R}(\vec x, y)).
  \end{align*}
  Kuantifikasi eksistensial mawa wates uga bisa ditegesi
  nganggo $\fn{max}$. Utawa, bisa ditegesi saka
  kuantifikasi universal mawa wates, nganggo ekuivalensi
  $\bexists{z < y}{R(\vec x, z)}
  \liff \lnot \bforall{z < y}{\lnot R(\vec x, z)}$.
  Gatekna yen, umpamane, kuantor mawa wates kanthi
  wujud $\bexists{x \leq y}{\dots
  x\dots}$ ekuivalen karo
  $\bexists{x < y+1}{\dots x \dots}$.
\end{proof}

\begin{prob}
  Tuduhna yen relasi telung panggonan $x \equiv y \mod n$
  (kongruensi modulo~$n$) iku rekursif primitif. Modulus biasane positif; yen modulus nol uga kalebu minangka input, kongruensi modulo nol ing gladhen iki ditegesi minangka kesamaan rong input.
\end{prob}

Fungsi rekursif primitif liyane kang migunani yaiku
fungsi kondisional, $\fn{cond}(x,y,z)$, kang ditegesi kanthi
\begin{align*}
  \fn{cond}(x,y,z) & = \begin{cases}
  y & \text{yen $x = 0$} \\
  z & \text{yen ora}.
\end{cases}
\intertext{Fungsi iki ditegesi kanthi rekursif liwat}
\fn{cond}(0,y,z) & = y,\\
\fn{cond}(x+1,y,z) & = z.
\end{align*}
Iki bisa digunakake kanggo mbenerake definisi fungsi
rekursif primitif miturut kasus saka relasi rekursif primitif:

\begin{prop}
Yen $g_0(\vec x)$, \dots,~$g_m(\vec x)$ iku fungsi rekursif
primitif, lan $R_0(\vec
x)$, \dots, $R_{m-1}(\vec x)$ iku relasi rekursif primitif,
mula fungsi $f$ kang ditegesi kanthi
\[
f(\vec x) = \begin{cases}
    g_0(\vec x) & \text{yen $R_0(\vec{x})$} \\
    g_1(\vec x) & \text{yen $R_1(\vec{x})$ lan dudu $R_0(\vec{x})$} \\
    \vdots & \\
    g_{m-1}(\vec x) & \text{yen $R_{m-1}(\vec{x})$ lan ora ana
      syarat sadurunge kang bener}
    \\
    g_m(\vec x) & \mbox{yen ora}
\end{cases}
\]
uga rekursif primitif.
\end{prop}

\begin{proof}
  Yen $m = 1$, iki mung fungsi kang ditegesi kanthi
  \[
  f(\vec x) = \fn{cond}(\Char{\lnot R_0}(\vec x),g_0(\vec x),g_1(\vec
  x)).
  \]
  Kanggo $m$ luwih gedhe tinimbang $1$, kita cukup
  nyusun komposisi definisi kanthi wujud iki.
\end{proof}
\end{document}
